\subsection{里斯自同态的谱特征刻画}

设 E 是复巴拿赫空间，u 是 E 的自同态。回忆我们用 \(\mathrm{Sp}(u)\) 表示 u 相对于含幺巴拿赫代数 \(\mathcal{L}(\mathrm{E})\) 的谱（参见 I, p. 127 的 § 7）。

假设 0 是 \(\mathrm{Sp}(u)\) 的一个孤立点；回忆这时我们用 \(e_0(u)\) 表示与 \(u\) 以及开闭子集 \(\{0\}\)（它是 \(u\) 的谱的子集）相伴的谱投影（参见 I, p. 131 的第 3 目）。我们有 \(e_0(u) = f(u)\)，它对每个全纯函数芽 \(f\)（在 \(\mathrm{Sp}(u)\) 的邻域中定义、在 0 的一个邻域中等于 1 且在 \(\mathrm{Sp}(u) \setminus \{0\}\) 的一个邻域中为零）成立。

通过过渡到子空间，自同态 \(u\) 在 \(e_0(u)\) 的像上诱导出一个拟幂零自同态（其谱约化为 \(\{0\}\)），并在 \(e_0(u)\) 的核上诱导出一个自同构（其谱为 \(\mathrm{Sp}(u) \setminus \{0\}\)）（loc. cit.）。

命题 14. --- 设 E 是复巴拿赫空间。为使 \(\mathcal{L}(\mathrm{E})\) 的元素 u 是 E 的里斯自同态，必须且只需下列两个互斥条件之一被满足：

\begin{enumerate}
\def\labelenumi{(\roman{enumi})}
\item
  E 的自同态 u 是 E 的自同构；
\item
  点 0 是 \(\operatorname{Sp}(u)\) 的孤立点，且投影 \(e_0(u)\) 有有限秩。
\end{enumerate}

当条件 (ii) 满足时，\(e_{0}(u)\) 的像是 u 的幂零空间，而 \(e_{0}(u)\) 的核是 u 的余幂零空间。

E 的每个自同构都是 E 的里斯自同态（III, p. 47, 注记 1）。如果 u 满足条件 (ii)，\(e_{0}(u)\) 的核 F 是 E 的闭向量子空间，有有限余维，且在 u 下稳定。设 \(u_{F}\) 是由 u 导出的 F 的自同态。它的谱包含在 \(\mathbf{C}\setminus\{0\}\) 中（I, p. 131 的第 3 目），因而 \(u_{F}\) 是 F 的自同构。由此得出 u 是 E 的里斯自同态（III, p. 48, prop. 8）。

反之，设 \(u\) 是 E 的里斯自同态。用 N 表示它的幂零空间，用 I 表示它的余幂零空间。由 III, p. 47 的 prop. 6，空间 E 是 N 与 I 的拓扑直和，且映射 \(u\) 通过过渡到子空间定义了 N 的一个幂零自同态 \(u_{\mathrm{N}}\) 与 I 的一个自同构 \(u_{\mathrm{I}}\)。特别地，我们有 \(\mathrm{Sp}(u_{\mathrm{N}}) \subset \{0\}\) 与 \(0 \not\in \mathrm{Sp}(u_{\mathrm{I}})\)。如果 N 是零空间，那么 I = E 且 \(u\) 是 E 的自同构。否则，0 是 \(\mathrm{Sp}(u)\) 的孤立点，且 \(e_0(u)\) 是以 N 为像、以 I 为核的投影（I, p. 129 的 prop. 2）；特别地，它有有限秩。

注记. --- 设 E 是实巴拿赫空间，u 是 E 的自同态。假设 0 是 u 的复谱的一个孤立点，也就是说，是自同态 \(1 \otimes u\)（它作用于 E 的完备可赋范复化空间 \(\mathrm{E}_{(\mathbf{C})}\)）的谱的一个孤立点（I, p. 85, 第 13 目）。子集 \(\{0\}\)（它属于 \(u\) 的复谱）是开闭的，且在共轭下不变；用 \(e_0(u) \in \mathcal{L}(\mathrm{E})\) 表示与之相伴的谱投影（I, p. 85 的第 13 目）。命题 14 在此情形下作必要的修改后成立。

